\documentclass{article}
\usepackage[utf8]{inputenc}
\usepackage{hyperref}

\title{Cahier des Charges \\ \textit{An Era of the Seas}}
\author{Dream Chasers}
\date{Octobre 2024}

\begin{document}

\maketitle










\newpage
\section*{Introduction}

\subsection*{Sommaire}
\tableofcontents










\newpage
\section{Origine et nature du projet}
\subsection{Origin}\hfill

Le projet du jeu An Era of the Seas est originaire d’un intérêt partage entre les membres de l’équipe dans les jeux vidéo et le développement de logiciels. Cela implique de combiner les connaissances en informatique et de développement de jeux vidéo pour créer une expérience captivante. Ce projet est une opportunité pour mettre en pratique les connaissances dans la création d’un produit concret.

Développer un jeu est toute une opportunité d'acquérir des compétences dans différents domaines, de la programmation à la conception graphique, sans oublier la production sonore. Chaque membre peut contribuer avec ses compétences uniques, ce qui enrichit l'expérience collective.

L’essor de l’industrie du jeu vidéo et l’importance croissante de son importance dans différents domaines tels que l’éducation et le divertissement ont également influencé la décision de créer un jeu. C’est l’occasion de s’immerger dans le secteur dynamique et en constante évolution du jeu vidéo.





\subsection{Information du jeu} 
\subsubsection{Début et histoire du personnage principal}\hfill

An Era of the Seas est un jeu dans lequel l’environnement maritime est crucial. Les joueurs incarneront le fils d’un grand et puissant marchant qui avait beaucoup d’influence et de contrôle sur les mers. Un jour il reçut une lettre lui informant que son père a péri durant une expédition mystérieuse. Il décide alors de trouver ce qui lui est arrivé en débutant son aventure en tant que capitaine, il souhaite également restaurer la gloire de son père. Il découvrira la diversité des iles et commence son voyage pour, non seulement devenir un grand héro comme son père l’était, mais aussi trouver ses propres réponses.



\subsubsection{Objectifs du joueur}\hfill

Les joueurs pourront naviguer et découvrir les iles de Shattered Seas. Pour commencer cette aventure, le capitaine n’aura seulement une barque en bois qu’il a trouvé sur la plage. Il devra explorer les 8 différentes iles. Les joueurs commenceront sur la première ile, ou ils découvriront les caractéristiques principales, les fonctionnalités et l’histoire du fils du grand marchand. Ils y découvriront le chemin pour accomplir l’objectif final du capitaine. Le capitaine devra voyager et atteindre chaque ile pour retrouver les trésors dispersés de son père pour pouvoir avoir le contrôle de Shattered Seas. Après avoir explorer les 7 premières iles, le capitaine pourra enfin terminer son aventure sur la dernière ile ou les difficultés et les défis restent à venir.



\subsubsection{Fonctionnalités du jeu}\hfill

Ce je présente des fonctionnalités uniques liées à son environnement

\paragraph{Quêtes principale}\hfill

Ce jeu est basé sur une histoire. Le joueur vivra l'histoire du capitaine et sera immergé dans l'univers maritime de ce jeu. Grâce à la quête principale, le joueur peut ensuite suivre l'histoire principale.

\paragraph{Quêtes secondaires}\hfill

La quête principale est l'histoire principale du capitaine, mais il y a aussi des petites histoires secondaires qui s'est déroulée au cours de son aventure à travers les mers et les îles. Ces quêtes peuvent aider le joueur à comprendre certaines fonctionnalités et offrent de généreuses récompenses. Ces quêtes consistent à faire de tâches, des livraisons, des combats contre certains ennemis spécifiques ou d'aider un résident de Shattered Seas.

\paragraph{Système de navigation immersive}\hfill

Le bateau dépendra du vent ou des rames pour naviguer en fonction du niveau du bateau et de l'équipage. Cela signifie naviguer avec des rames pour les bateaux ayant un niveau faible, puis suivre la direction du vent, régler les voiles et utiliser le gouvernail pour les bateaux ayant un niveau supérieur. Pour que le capitaine puisse se localiser et trouver un chemin pour qu’il puisse se rendre à un endroit en mer, il utilisera une carte et une boussole.

\paragraph{Entretien et améliorations du bateau/navire}\hfill

Le bateau du joueur peut être amélioré en fonction de l’avancement de la quête de l’histoire principale. Lorsque le bateau est cassé, le joueur devra le réparer en utilisant des matériaux pour pouvoir à nouveau naviguer sur les mers. Il pourra acheter la base du bateau auprès d'un vendeur et l'améliorer avec certains matériaux.

\paragraph{Membre d’équipage}\hfill

Puisque le joueur est un capitaine naviguant sur les mers, il dispose de son propre équipage. Cet équipage est composé de différents membres aux tâches multiples qui ont été soigneusement embauchés dans les villes et des quêtes dans le but d'aider le joueur dans son voyage. Chaque équipage possède des statistiques spécifiques. Il est conseillé de les attribuer aux bonnes tâches au profit du parcours du capitaine. Par exemple, certains membres de l’équipage savent mieux gérer les rames et d’autres sont meilleurs en exploration.

\paragraph{Armes}\hfill

Pour pouvoir lutter contre les ennemis, le joueur a la possibilité d'utiliser des épées et des fusils. Les armes peuvent être obtenues dans des coffres au trésor, mais celles-ci sont courantes et faibles. Il est recommandé de les acheter chez un marchand dans les villes pour obtenir des armes avec des statistiques plus correctes et infliger plus de dégâts.

\paragraph{Stigmata (Améliorations du joueur)}\hfill

Les Stigmata (Stigma au singulier) sont des marques à la surface de la main, créées ou inscrites sur la main avec l'aide d'un mage situé dans les villes. Ces marques valorisent les capacités du joueur (vitesse, attaque...). Les Stigmata peuvent être obtenus en les collectant dans des coffres au trésor ou en les achetant auprès du mage. Ces mages ont consacré beaucoup de dévouement et d’efforts pour créer des Stigmata. Par conséquent, les Stigmata vendus par les mages sont plus efficaces et ont plus de puissance et d'effet sur le capitaine. Cela signifie que le joueur obtient des statistiques plus correctes en achetant des Stigmata à un mage au lieu de les inscrire sur sa main avec l'aide du mage. Cette qualité n'est évidemment pas sans impact sur le nombre de pièces du capitaine, celles-ci coûtent bien plus cher que celles que l'on trouve dans la nature.

\paragraph{Coffre aux trésors}\hfill

Pour obtenir des ressources et des pièces, des coffres au trésor peuvent être trouvés sur l'île et dans des endroits gardés par les ennemis. Le joueur peut obtenir des pièces de monnaie, des armes, des stigmates et des objets pour améliorer le bateau, des armes et des Stigmata. Cela aide le capitaine à établir une certaine richesse et, comme l'argent signifie le pouvoir sur Shattered Seas, il pourrait avoir plus d'influence sur les mers.

\paragraph{Villes}\hfill

Chaque île a sa propre ville. Il y a plusieurs commerçants dans ces villes : 
\begin{itemize}
    \item \textbf{Les Fabricants de bateaux et de navires :} Ils vendent leur beau travail de fabrication de bateaux et de navires imposants, selon l'île, ces concessionnaires proposent de meilleures coques et améliorations de navires.
    \item \textbf{Les Fabricants d'armes :} Ils vendent des épées et des fusils efficaces pour affronter les ennemis.
    \item \textbf{Les Mages :} Ils peuvent utiliser leur magie mystique pour inscrire une Stigma que vous possédez ou que vous avez achetée. Ils vendent aussi leurs propres Stigmata qui sont plus puissants.
    \item \textbf{Les Vendeurs de nourriture :} ils vendent de la nourriture délicieuse aux habitants et aux aventuriers dans l'optique de ne pas mourir de faim et d'avoir suffisamment d'énergie.
\end{itemize}

\paragraph{Combat (PvE)}\hfill

Dans cette aventure, le chemin du capitaine n’est évidemment pas paisible, il rencontrera divers ennemis aussi bien sur les îles qu’à travers les mers. Pour défendre et attaquer les ennemis, les armes, Stigmata du capitaine et l’équipage seront utiles sur le terrain. Comme sur les mers, le bateau saura attaquer, et les membres d'équipage sont essentiels en la matière. Les statistiques des armes, des stigmates, de l’équipage et du bateau sont donc essentielles pour infliger des dégâts et résister aux attaques des ennemis.

\paragraph{Multijoueur}\hfill

Ce jeu a un mode multijoueur. Dans ce mode, le joueur pourra choisir une classe et jouer selon cette classe. Notez que dans ce mode, le joueur n'a aucun avantage en fonction de l'avancement et des objets du mode Solo.

\paragraph{Pèche}\hfill

Une mécanique utile et amusante grâce à laquelle le joueur peut obtenir des poissons utiles pour se nourrir et des trésors comme des pièces et des ressources. Cela sera possible grâce à une canne à pêche, remise au joueur en début de partie et sera autorisée dans des plans d'eau assez grands comme les lacs et les océans.





\subsection{Outils}
\subsubsection{Outils de gestion d’équipe}\hfill

Pour pouvoir communiquer différents aspects du développement du jeu entre les membres de l’équipe, il est crucial de disposer d’un outil dédié à cet effet. Ainsi, un serveur Discord est particulièrement utile. Cela permet à l'équipe de communiquer facilement et les discussions peuvent être organisées en différents canaux en fonction des sujets de développement spécifiques.
\\

De plus, la gestion des tâches est essentielle pour maintenir le projet sur la bonne voie et suivre chaque avancement. Une liste de choses à faire peut être très utile dans ce domaine. Il peut fournir un aperçu clair des tâches à effectuer. Microsoft To Do est une excellente solution pour cela. Les tâches peuvent être attribuées à un membre de l'équipe et la date d'échéance peut être définie pour une tâche. Cela permet au membre de l'équipe d'être conscient de ce qui doit être fait, des délais et de maintenir sa productivité.



\subsubsection{Outils de développement du jeu}\hfill

Les Dream Chasers utilisent GitHub pour pouvoir travailler correctement ensemble sur le projet. Il permet à l’équipe de suivre les modifications apportées au code du jeu, garantissant ainsi que tout le monde travaille sur la dernière version sans écraser le travail des autres.\\

Le jeu est codé en C\# (Csharp) pour plusieurs raisons. Par exemple, C\# est un langage de programmation orienté objet, une caractéristique utile pour les jeux et C\# est également très souvent utilisée dans les moteurs de jeux comme Unity.\\

Pour l’IDE, Rider de JetBrains car il présente divers avantages. Il est très puissant spécifiquement pour le développement .NET (dotnet). Ici, l'équipe utilise dotnet core 7 pour développer le jeu. L'outil de débogage de cet IDE est très puissant, il facilite l'identification des problèmes dans le code grâce au débogage étape par étape et à la surveillance des variables. Rider fournit un éditeur Unity, et il est également plus rapide que certains autres IDE comme Visual Studio.\\

Le moteur de jeu utilisé est Unity comme mentionné précédemment pour de nombreuses raisons. Il offre une interface utilisateur accessible qui facilite la gestion des projets par les développeurs. Avec Unity, le jeu peut être pris en charge sur plusieurs plateformes. Ainsi le jeu sera accessible même si le joueur utilise Windows, macOS ou Linux. Il fournit une variété d'outils et de fonctionnalités pour l'environnement 3D du jeu, sans mentionner qu'un moteur physique est inclus.\\

Pour stocker les données du joueur, les données sont stockées en ligne sur Firebase Realtime Database de Google. Cela offre une synchronisation en temps réel. L'interface utilisateur facilite son utilisation par les développeurs. Les données sont stockées au format JSON et Firebase permettent de manipuler (envoyer et lire des données) facilement. Le joueur doit se connecter pour jouer et l'authentification Firebase rend la chose facile pour les développeurs et est sécurisée avec une gestion des utilisateurs très utile.



\subsubsection{Budget}










\newpage
\section{Objet d’étude}
\subsection{Intérêts et objectifs}\hfill
\subsection{Bénéfices commun}\hfill
\subsection{Bénéfices individuels}\hfill










\newpage
\section{Sid Meier’s Pirates}
\subsection{Sid Meier’s Pirates}
\subsubsection{Point forts}
\paragraph{Cadre historique et immersion}\hfill

\paragraph{Choix et liberté du joueur}\hfill

\paragraph{Possibilité de rejouer}\hfill



\subsubsection{Fonctionnalités}
\paragraph{Factions et diplomatie}\hfill

\paragraph{Combat d’épée et duels}\hfill

\paragraph{Romance et réputation}\hfill

\paragraph{Quêtes secondaire et Histoire personnel}\hfill





\subsection{Pirates of the Burning Sea}
\subsubsection{Points forts}
\paragraph{Combat naval}\hfill

\paragraph{Vie de pirate}\hfill



\subsubsection{Fonctionnalités}
\paragraph{Conquêtes de ports et contrôle de territoires}

\paragraph{Guildes et alliances}





\subsection{Sea of Thieves}
\subsubsection{Points forts}
\paragraph{Gameplay social et coopératif}\hfill

\paragraph{Contenu évolutif et élargi}\hfill

\paragraph{Liberté}\hfill



\subsubsection{Fonctionnalités}
\paragraph{World Events}\hfill

\paragraph{Environnement interactif}\hfill

\paragraph{Accent sur le choix des joueurs}\hfill










\newpage
\section{Entreprise}
\subsection{Nom}\hfill



\subsection{Spécificités}\hfill



\subsection{Domaine d’activités}\hfill










\newpage
\section{Equipe}
\subsection{Jérôme}\hfill



\subsection{Achref}\hfill



\subsection{Ethan}\hfill



\subsection{Joseph}\hfill



\subsection{Matteo}\hfill










\newpage
\section{Distribution des taches}










\newpage
\section{Avancement et planification }


\end{document}